\appsection{Identification from Choice Probabilities under EUT}
\label{app:eut_choice_probs}
%==============================================================================

The quantity-space prior specification of Section~\ref{sec:framework}
requires an injective mapping from parameters to predictive quantities,
and the Jacobian pullback must faithfully recover the intended prior on
the feasible image. As noted in Section~\ref{sec:discussion}, certainty
equivalents require inverting the utility function, which complicates the
injectivity proof for richer utility families. Choice probabilities
avoid this complication. This appendix proves that, under Expected
Utility Theory with a logit choice rule, the Jacobian pullback from
choice probabilities recovers the intended prior for any analytic
utility family, provided the lottery design satisfies a rank condition
that is easy to verify.

The key observation is that the map from $p$ utility parameters to $p$
pairwise choice probabilities factors as
$G = \Lambda \circ (C\phi + \mathbf{d})$, where $\phi$ evaluates utility
at interior outcomes, $C$ is a contrast matrix determined by the lottery
design, and $\Lambda$ is the coordinatewise logistic link. Since
$\Lambda$ is a bijection with strictly positive derivative, all
identification content lies in the product $C\,D\phi$. The results here
apply to EUT specifically; extending to non-expected utility functionals
such as RDU, where the probability weighting layer breaks the affine
structure of $C\phi + \mathbf{d}$, is left for future work.


\subsection{Setup}

Let $\Theta \subseteq \mathbb{R}^p$ be open and bounded, and let
$\{u(\cdot;\,\boldsymbol{\theta})\}_{\boldsymbol{\theta} \in \Theta}$
be a family of utility functions on $[x_{\min}, x_{\max}]$ such that for
each $\boldsymbol{\theta}$, $x \mapsto u(x;\,\boldsymbol{\theta})$ is
strictly increasing with $u(x_{\min};\,\boldsymbol{\theta}) = 0$ and
$u(x_{\max};\,\boldsymbol{\theta}) = 1$. Assume $u$ is analytic in
$\boldsymbol{\theta}$ and $C^1$ in $(x,\boldsymbol{\theta})$.

We observe $R = p$ pairwise comparisons. Across all comparisons, let
\begin{equation}
  x_{\min} = y_0 < m_1 < \cdots < m_M < y_{M+1} = x_{\max}
\end{equation}
be the common support obtained by taking the distinct outcomes that
appear anywhere in the design and adjoining the normalization
endpoints if necessary. For comparison $r$, let
$w_{r\ell} \coloneqq \Pr_{A_r}(y_\ell) - \Pr_{B_r}(y_\ell)$ be the
signed probability contrast at outcome $y_\ell$. Under EUT, the
expected-utility difference is
\begin{equation}\label{eq:delta_common_app}
  \Delta_r(\boldsymbol{\theta})
  =
  \sum_{\ell=0}^{M+1} w_{r\ell}\,u(y_\ell;\,\boldsymbol{\theta}).
\end{equation}

Define the \textit{evaluation map}
\begin{equation}\label{eq:phi_def_app}
  \phi(\boldsymbol{\theta})
  =
  \bigl(u(m_1;\,\boldsymbol{\theta}), \ldots,
        u(m_M;\,\boldsymbol{\theta})\bigr)
  \in (0,1)^M,
\end{equation}
which records utility at the $M$ interior outcomes. Since
$u(x_{\min};\,\boldsymbol{\theta}) = 0$ kills the $\ell = 0$ term and
$u(x_{\max};\,\boldsymbol{\theta}) = 1$ reduces the $\ell = M{+}1$ term
to the contrast weight alone, \eqref{eq:delta_common_app} simplifies to
\begin{equation}\label{eq:Delta_linear_app}
  \boldsymbol{\Delta}(\boldsymbol{\theta})
  =
  C\,\phi(\boldsymbol{\theta}) + \mathbf{d},
\end{equation}
where $C = (w_{ri})$ is the $p \times M$ \textit{contrast matrix} and
$d_r = w_{r,\,M+1}$. The observed choice probability is
$\pi_r(\boldsymbol{\theta}) = \Lambda(\Delta_r(\boldsymbol{\theta}))$
with $\Lambda(t) = e^t/(1+e^t)$, giving the
\textit{choice-probability map}
\begin{equation}\label{eq:factorization_app}
  G(\boldsymbol{\theta})
  =
  \bigl(\pi_1(\boldsymbol{\theta}), \ldots,
        \pi_p(\boldsymbol{\theta})\bigr)
  =
  \Lambda\!\bigl(C\,\phi(\boldsymbol{\theta}) + \mathbf{d}\bigr).
\end{equation}

Normalization matters here in a way it does not for certainty
equivalents. For certainty equivalents, positive affine transformations
of utility are innocuous. For logit probabilities they are not: additive
shifts cancel in $\Delta_r$, but rescalings do not. Without endpoint
normalization, the utility scale would be confounded with the logit noise
scale.


\subsection{Identification}

\begin{theorem}[Identification from $p$ choice probabilities]
\label{thm:choice_prob_id}
Let $G \colon \Theta \to (0,1)^p$ be the choice-probability map
\eqref{eq:factorization_app}.
\begin{enumerate}
  \item \textit{Jacobian.} $J_G(\boldsymbol{\theta}) =
    D_\Lambda(\boldsymbol{\theta})\,C\,D\phi(\boldsymbol{\theta})$, where
    $D_\Lambda = \diag(\pi_r(1-\pi_r))_{r=1}^p$, which is a
    positive diagonal matrix since each $\pi_r \in (0,1)$. Hence
    \begin{equation}\label{eq:rank_equiv_app}
      \operatorname{rank}(J_G(\boldsymbol{\theta}))
      = \operatorname{rank}\!\bigl(C\,D\phi(\boldsymbol{\theta})\bigr).
    \end{equation}
  \item \textit{Singular set.} At some parameter values the
    Jacobian may lose rank, so that $G$ is not locally injective
    there. Call these the \textit{singular set}
    $S = \{\boldsymbol{\theta} :
      \operatorname{rank}(J_G) < p\}$. If
    $\operatorname{rank}(C\,D\phi(\boldsymbol{\theta}_0)) = p$ at even
    one point $\boldsymbol{\theta}_0 \in \Theta$, then $S$ has Lebesgue
    measure zero and $\Theta \setminus S$ is open and dense. In other
    words, rank failure is the exception, not the rule: it can occur,
    but only on a negligible set of parameter values.
  \item \textit{Local identification.} At every
    $\boldsymbol{\theta} \in \Theta \setminus S$ the Jacobian has full
    rank, so $G$ is locally injective: nearby parameter values produce
    distinct choice probabilities.
  \item \textit{Bayesian pullback.} Let $q$ be a prior density on
    $(0,1)^p$. If, for almost every
    $\boldsymbol{\pi} \in G(\Theta)$, the equation
    $G(\boldsymbol{\theta}) = \boldsymbol{\pi}$ has exactly $m$
    solutions, then the pullback
    \begin{equation}
      \tilde q(\boldsymbol{\theta})
      \propto
      q(G(\boldsymbol{\theta}))\,|\det J_G(\boldsymbol{\theta})|
    \end{equation}
    induces the density $m\,q(\boldsymbol{\pi})$ on $G(\Theta)$.
    After normalization, this is exactly the intended prior restricted
    to the feasible image $G(\Theta)$. Any prior mass that $q$ assigns
    outside $G(\Theta)$ is irrelevant because no parameter value can
    reach it. When $m = 1$, no multiplicity correction is needed.
\end{enumerate}
\end{theorem}

\begin{proof}
\textit{Part~1.} Since $G = \Lambda \circ (C\phi + \mathbf{d})$ and
$\Lambda'(t) = \Lambda(t)(1-\Lambda(t)) > 0$, the chain rule gives
$J_G = D_\Lambda\,C\,D\phi$. Since $D_\Lambda$ is a positive
diagonal matrix, multiplying by it does not change rank.

\textit{Part~2.} Whether $J_G$ has full rank at $\boldsymbol{\theta}$
depends, by Part~1, on whether
$C\,D\phi(\boldsymbol{\theta})$ has full rank, which in the square case
($M = p$) amounts to whether
$\det(C\,D\phi(\boldsymbol{\theta}))$ is nonzero. Write
$f(\boldsymbol{\theta}) = \det(C\,D\phi(\boldsymbol{\theta}))$.
Because $u$ is analytic in $\boldsymbol{\theta}$, the function $f$
is analytic. Because there exists $\boldsymbol{\theta}_0$ such that
$f(\boldsymbol{\theta}_0) \neq 0$, the function $f$ is not
identically zero. A standard property of analytic functions is that
a nonzero analytic function can only vanish on a set of Lebesgue
measure zero \parencite{Krantz2002,Mityagin2015}. Therefore
$S = \{f = 0\}$ has measure zero. When $M > p$, the same argument
applies to any $p \times p$ subdeterminant of $C\,D\phi$ that is
nonzero at $\boldsymbol{\theta}_0$, and $S$ is contained in its
zero set.

\textit{Part~3.} At each
$\boldsymbol{\theta} \in \Theta \setminus S$,
$J_G(\boldsymbol{\theta})$ has full rank $p$. By the inverse function
theorem, $G$ is locally injective at such $\boldsymbol{\theta}$: no
nearby parameter value produces the same choice probabilities.

\textit{Part~4.} Since $\Lambda$ is strictly monotone,
$G(\boldsymbol{\theta}_1) = G(\boldsymbol{\theta}_2)$ if and only if
$C\phi(\boldsymbol{\theta}_1) + \mathbf{d} =
 C\phi(\boldsymbol{\theta}_2) + \mathbf{d}$.
So $G$ has the same number of preimages at each point as the
utility-difference map $C\phi + \mathbf{d}$. The multivariate
change-of-variables formula then gives, for every measurable
$B \subseteq G(\Theta)$,
\begin{equation}
  \int_{G^{-1}(B)}
    q(G(\boldsymbol{\theta}))\,|\det J_G(\boldsymbol{\theta})|\,
    d\boldsymbol{\theta}
  =
  \int_B q(\boldsymbol{\pi})\,\#G^{-1}(\boldsymbol{\pi})\,
    d\boldsymbol{\pi}
  =
  m \int_B q(\boldsymbol{\pi})\,d\boldsymbol{\pi}.
  \qedhere
\end{equation}
\end{proof}

Identification requires $M \geq p$ interior outcomes and
$\operatorname{rank}(C) \geq p$. If any lottery pair has $A_r$
first-order stochastically dominating $B_r$, the corresponding
$\Delta_r$ is nonnegative for every increasing utility, so that
comparison fixes the sign of the EU gap. Such comparisons can still
be informative about utility parameters through the magnitude of
$\Delta_r$, so non-FOSD is not necessary for identification. It is,
however, often a useful design feature because non-FOSD pairs allow
the sign of the EU gap itself to vary with preferences. In any case,
one still needs $C\,D\phi$ to have rank $p$.


\subsection{Sufficient Designs}

The theorem applies to any lottery design whose contrast matrix $C$
and evaluation Jacobian $D\phi$ jointly produce rank $p$. Two examples
illustrate this.

\textit{Sure amount versus endpoint lottery.} Fix interior outcomes
$m_1,\ldots,m_p \in (x_{\min},x_{\max})$ and
$\lambda_1,\ldots,\lambda_p \in (0,1)$. For each $i$, compare a sure
amount $m_i$ against a lottery paying $x_{\max}$ with probability
$\lambda_i$ and $x_{\min}$ otherwise. Normalization gives
$\operatorname{EU}(B_i) = \lambda_i$, so
$\Delta_i(\boldsymbol{\theta}) =
 u(m_i;\,\boldsymbol{\theta}) - \lambda_i$.
The only interior outcome in pair $i$ is $m_i$, with contrast weight
$+1$, so $C = I_p$ and the Jacobian simplifies to
\begin{equation}\label{eq:binary_jac_app}
  J_G(\boldsymbol{\theta})
  =
  \diag\!\bigl(\pi_i(1-\pi_i)\bigr)_{i=1}^p
  \cdot
  \left(
    \frac{\partial u(m_i;\,\boldsymbol{\theta})}{\partial \theta_j}
  \right)_{\!i,j}.
\end{equation}
With this design, identification reduces entirely to the evaluation
map: the utility family must be such that the matrix of partial
derivatives $(\partial u(m_i;\,\boldsymbol{\theta})/\partial
\theta_j)_{i,j}$ has full rank at some $\boldsymbol{\theta}_0$.

\textit{Risk-apportionment tasks.} A risk-apportionment (RA) task of
order $k$ is a binary lottery comparison in which the first $k-1$
moments of the two lotteries are matched and the $k$th moments differ
\parencite{Eeckhoudt2006}. For example, RA2 (risk aversion) matches
means but differs in variance, RA3 (prudence) matches means and
variances but differs in skewness, and so on. Each RA task contributes
a row to $C$. Unlike the binary design, $C$ is no longer the identity,
so the design-side rank condition is not automatic. However, it is
straightforward to verify numerically for any given set of RA tasks.
RA tasks are a natural choice because each task targets a different
order of the utility function, so for utility families whose parameters
control these derivatives somewhat independently, the rows of
$C\,D\phi$ will generically point in different directions.

\subsection{Noise Parameter and Sequential Composition}

The results above identify $p$ utility parameters from $p$ choice
probabilities. In practice, the choice model also includes a noise
parameter $\lambda$ that scales the logit link. This can be handled by
sequential composition: first identify the $p$ utility parameters using
$p$ lottery pairs as above, then identify $\lambda$ with one additional
pair whose choice probability varies with $\lambda$ given the utility
parameters. The $(p+1)$th pair contributes one additional row to the
system, and the same Jacobian pullback applies to the extended
$(p+1)$-dimensional map.
